\section{Founding populations and catastrophe}
\label{sec:small}

Criticality controls long-time behaviour, but a small population experiences most
of its immediate extinction risk before that asymptotic distinction is visible.
At \(p=\tfrac12\), for example, the extinction probabilities in the first four
generations are
\begin{equation}
  u_1=0.500,
  \qquad u_2=0.625,
  \qquad u_3=0.695,
  \qquad u_4=0.742
  \quad\text{(three significant figures)}.
  \label{eq:early-extinction-values}
\end{equation}
Survival to generation two has probability \(1-u_2=0.375\): the founder must
divide, and at least one of its two daughters must then divide.  Small movements
of \(p\) around one half have little time to accumulate over these first updates,
so early extinction remains substantial on both sides of the critical point.

\subsection{A cohort of \texorpdfstring{\(k\)}{k} founders}
\label{sec:cohort}

Suppose \(Z_0=k\).  The descendants of the \(k\) founders form independent
lineages.  If \(u_n\) is the extinction probability by generation \(n\) from one
founder, all \(k\) lineages are extinct with probability \(u_n^k\), and therefore
\begin{equation}
  S_n^{(k)}=\Prb(Z_n>0\mid Z_0=k)=1-u_n^k.
  \label{eq:Snk}
\end{equation}
This formula separates the effect of founding count from the dynamics of one
lineage.  Taking the long-time limit gives
\begin{equation}
  S_\infty^{(k)}=
  \begin{cases}
    0,&0\leq p\leq\tfrac12,\\[3pt]
    \displaystyle
    1-\left(\frac{1-p}{p}\right)^k,
    &\tfrac12<p\leq1.
  \end{cases}
  \label{eq:cohort-ultimate-survival}
\end{equation}
A larger cohort cannot prevent eventual extinction at or below criticality.  In
the supercritical range it instead supplies repeated independent opportunities
for one lineage to establish.

The finite-time effect is also explicit.  For fixed \(k\), expansion of
\(1-(1-S_n)^k\) at small \(S_n\) gives
\begin{equation}
  S_n^{(k)}\sim
  \begin{cases}
    kA(p)(2p)^n,&p<\tfrac12,\\[3pt]
    2k/n,&p=\tfrac12.
  \end{cases}
  \label{eq:cohort-late-asymptotics}
\end{equation}
The cohort multiplies the late survival probability by \(k\), to leading order,
but it does not change the subcritical Yaglom scale.  Indeed,
\begin{equation}
  \E[Z_n\mid Z_n>0,Z_0=k]
  =\frac{k(2p)^n}{S_n^{(k)}}
  \longrightarrow\frac1{A(p)}.
  \label{eq:cohort-yaglom-mean}
\end{equation}
At late times, conditioning selects a surviving lineage; the number of failed
founders no longer affects the conditional population size.

Conditioning on survival can also bias the history that is reconstructed from
successful populations.  Budd and Mann's analysis of phylogenetic clades shows
that clades retained because they survive to the present are enriched for high
early origination rates even under homogeneous birth--death dynamics
\cite{Budd2018}.  This ``push of the past'' is related to
\eqref{eq:cohort-yaglom-mean}: selection acts on realised trajectories, so an
apparently exceptional early expansion need not imply a change in the underlying
per-capita rates.  The parallel is a modelling warning rather than an assertion
that phylogenetic and within-compartment mechanisms are the same.

\subsection{A discrete catastrophe mechanism}
\label{sec:bdc}

The later compartment models add a second absorbing outcome: rupture or
catastrophe.  A minimal discrete construction is useful because it distinguishes
three quantities that are otherwise easily conflated: the hazard of catastrophe
at the next update, the probability of catastrophe at an exact update, and the
probability of remaining both unruptured and non-extinct.

Let each individual present at the beginning of a generation trigger catastrophe
independently with probability \(\kappa\in[0,1]\), and put
\(c=1-\kappa\).  Conditional on a population of size \(z\), the probability that
no individual triggers catastrophe is \(c^z\).  If catastrophe is avoided, each
individual then follows the binary offspring law \eqref{eq:pgfbinary}; catastrophe
terminates the process.

First take pure division, \(p=1\).  Conditional on reaching update \(j\), the
population has size \(2^j\), so the catastrophe hazard is
\begin{equation}
  h_j=1-c^{2^j}.
  \label{eq:pure-cat-hazard}
\end{equation}
Avoiding catastrophe for \(n\) complete updates requires every individual in
generations \(0,\ldots,n-1\) to avoid triggering it, and hence
\begin{equation}
  C_n=c^{\sum_{j=0}^{n-1}2^j}=c^{2^n-1}.
  \label{eq:bdcS}
\end{equation}
The unconditional probability that catastrophe first occurs at update \(j\) is
therefore
\begin{equation}
  \Prb(K=j)=C_jh_j
  =c^{2^j-1}\bigl(1-c^{2^j}\bigr).
  \label{eq:pure-cat-mass}
\end{equation}
The factor \(C_j\) is indispensable: \(h_j\) alone is a conditional hazard, not an
exact-time probability.

When death is restored, the census before an update is random.  The process can
still be described exactly by a weighted PGF.  For an underlying binary branching
tree starting from one individual, define
\begin{equation}
  F_n(s)=
  \E\!\left[
    s^{Z_n}c^{\sum_{j=0}^{n-1}Z_j}
  \right],
  \qquad F_0(s)=s.
  \label{eq:cat-weighted-pgf}
\end{equation}
Conditioning on the founder's reproduction gives the closed recursion
\begin{equation}
  F_{n+1}(s)
  =c\,\phi(F_n(s))
  =c\bigl[1-p+pF_n(s)^2\bigr].
  \label{eq:cat-weighted-recursion}
\end{equation}
Its evaluations have distinct meanings:
\begin{align}
  F_n(1)
  &=\Prb(\text{no catastrophe in the first }n\text{ updates}),
  \label{eq:cat-no-rupture}\\
  F_n(1)-F_n(0)
  &=\Prb(Z_n>0\ \text{and no catastrophe in the first }n\text{ updates}).
  \label{eq:cat-alive-unruptured}
\end{align}
For \(p=1\), \eqref{eq:cat-weighted-recursion} reduces to
\eqref{eq:bdcS}; for \(\kappa=0\), it reduces to the ordinary
Galton--Watson iteration.

A tempting approximation replaces the accumulated random census in
\eqref{eq:cat-weighted-pgf} by its mean:
\begin{equation}
  F_n(1)\approx
  c^{\sum_{j=0}^{n-1}\E[Z_j]}
  =c^{\sum_{j=0}^{n-1}(2p)^j}.
  \label{eq:bdcguess}
\end{equation}
For \(0<c\leq1\), the map \(x\mapsto c^x\) is convex, so Jensen's inequality
shows that the right-hand side is a lower bound for
\emph{non-catastrophe alone}.  It is not a lower bound for
\eqref{eq:cat-alive-unruptured}, because extinction then contributes to
\(F_n(1)\) but not to survival.  The exact recursion
\eqref{eq:cat-weighted-recursion} is cheap enough that the approximation is not
needed here; its value is to show why a mean census cannot generally be inserted
into a path-dependent catastrophe law without specifying the event being bounded.

Whether a process conditioned jointly on non-extinction and non-catastrophe has a
quasi-stationary limit is a separate spectral question.  No such limit follows
from the Galton--Watson result alone, because the factor \(c^{Z_j}\) reweights the
population paths.  Later catastrophe models define their absorbing states and
conditioning events explicitly before making that inference.

\FloatBarrier
